\documentclass[12pt]{article}
\usepackage{graphicx}
\usepackage{amsmath}
\usepackage{geometry}
\geometry{a4paper, margin=1in}
\usepackage{titlesec}
\usepackage{minted}

% Format sections to match the reference document style
\titleformat{\section}
  {\normalfont\Large\bfseries}{\thesection.}{1em}{}
\titleformat{\subsection}
  {\normalfont\large\bfseries}{\thesubsection}{1em}{}

\begin{document}

\begin{center}
    {\Large \textbf{Quantum Simulation Lab}} \\
    \vspace{0.8cm}
    {\Large Experiment II} \\
    \vspace{0.2cm}
    {\Large \textbf{Deutsch Algorithm}} \\
\end{center}

\vspace{0.5cm}
\noindent Name: Hrishikesh M Jayakrishnan \hfill ID: SC26D054
\vspace{1em}
\hrule

\vspace{1em}
\section{Aim}
\begin{enumerate}
    \item To implement and demonstrate the Deutsch algorithm using a quantum circuit.
    \item To verify the correctness of the Deutsch algorithm in determining whether a given function is constant or balanced using the Qiskit framework.
\end{enumerate}
\section{Resources used}
\begin{itemize}
    \item Python
    \item Qiskit 
    \item Jupyter Notebook environment
\end{itemize}

\section{Theoretical Background}
\subsection{Deutsch Algorithm}
The Deutsch algorithm is a quantum algorithm that determines whether a given function is constant or balanced. A function is constant if it returns the same value for all inputs, and balanced if it returns 0 for half of the inputs and 1 for the other half. A classical algorithm would require atleast two evaluations of the function to determine whether it is constant or balanced, while the Deutsch algorithm can determine this with a single evaluation using quantum parallelism.
\newpage
\subsection{Steps}
Below are the steps involved in implementing the Deutsch algorithm for a given function \( f: \{0, 1\} \rightarrow \{0, 1\} \):
\begin{enumerate}
    \item Initialize two qubits, the first qubit in the state \( |0\rangle \) and the second qubit in the state \( |1\rangle \).
    \item Apply a Hadamard gate to both qubits to create a superposition of states.
    \item Apply the oracle \( U_f \) corresponding to the function \( f \) to the two qubits. The oracle encodes the function's output into the phase of the quantum state.
    \item The oracle is defined as follows: \[ U_f |x\rangle|y\rangle = |x\rangle|y \oplus f(x)\rangle \]
    \item Apply a Hadamard gate to the first qubit to interfere the amplitudes and extract information about the function.
    \item Measure the first qubit. If the result is \( |0\rangle \), the function is constant; if the result is \( |1\rangle \), the function is balanced.
\end{enumerate}


\section{Quantum circuit}
The experiment is conducted following these steps in the Qiskit framework:
\begin{enumerate}
    \item Initialize the quantum circuit with two qubits and two classical bits.
    \item Apply a Hadamard gate to the first qubit.
    \item Apply a Pauli-X gate to the second qubit.
    \item Apply the oracle \( U_f \) to the two qubits.
    \item Apply a Hadamard gate to the first qubit.
    \item Measure the first qubit.
\end{enumerate}

\begin{figure}[h!]
    \centering
    \includegraphics[width=1.02\textwidth]{../deutsch_algorithm.png}
    \caption{Quantum circuit for the Deutsch algorithm}
\end{figure}
\newpage
\section{Measurement Results and Analysis}
The circuit is executed on a quantum simulator for a single shot, and the measurement results are obtained. The results indicate whether the function is constant or balanced based on the measurement of the first qubit.
\begin{figure}[h!]
    \centering
    % Remove the \rule line below and uncomment the \includegraphics line to insert your image
    \includegraphics[width=0.5\textwidth]{../constant_0.png}
    % \rule{0.5\textwidth}{0.3\textwidth} % This creates a placeholder box
    \caption{Measurement result for a constant function \( f(x) = 0 \)}
\end{figure}
\begin{figure}[h!]
    \centering
    % Remove the \rule line below and uncomment the \includegraphics line to insert your image
    \includegraphics[width=0.5\textwidth]{../balanced.png}
    % \rule{0.5\textwidth}{0.3\textwidth} % This creates a placeholder box
    \caption{Measurement result for a balanced function \( f(x) = x \)}
\end{figure}
\newpage
The measurement results show that the circuit correctly identifies the nature of the function based on the measurement of the first qubit. If the result is \( |0\rangle \), the function is constant; if the result is \( |1\rangle \), the function is balanced, in a single evaluation of the function, demonstrating the efficiency of the Deutsch algorithm over classical methods.

\section{Conclusion}
This experiment implements the Deutsch algorithm using a quantum circuit in the Qiskit framework. The results obtained from the quantum simulator confirm that the algorithm can determine whether a function is constant or balanced with a single evaluation

\newpage
\section{Appendix}
\subsection{Python Code}
\begin{minted}[fontsize=\small]{python}
from qiskit import QuantumCircuit, QuantumRegister, ClassicalRegister, transpile
from qiskit_aer import AerSimulator
from qiskit.visualization import plot_histogram
import matplotlib.pyplot as plt

def create_oracle(function_type):
    oracle = QuantumCircuit(2, name="Oracle")
    if function_type == "constant_0":
        pass
    elif function_type == "constant_1":
        oracle.x(1)
    elif function_type == "balanced_same_input":
        oracle.cx(0, 1)
    elif function_type == "balanced_not":
        oracle.cx(0, 1)
        oracle.x(1)
    return oracle

x = QuantumRegister(1, "x")
y = QuantumRegister(1, "y")

cr = ClassicalRegister(1, "c")

qc = QuantumCircuit(x, y, cr)
fig = qc.draw("mpl")
plt.show()

qc.h(x)
qc.x(y)
qc.h(y)
fig = qc.draw("mpl")
plt.show()

oracle = create_oracle("balanced_same_input")
qc.append(oracle, [0, 1])

fig = qc.draw('mpl')
plt.show()

qc.barrier()
qc.h(x)
qc.measure(x, cr)
fig = qc.draw("mpl")
plt.show()

aer_sim = AerSimulator()
transpiled_qc = transpile(qc, aer_sim)
job = aer_sim.run(transpiled_qc, shots=1)
result = job.result()
counts = result.get_counts(transpiled_qc)
fig = plot_histogram(counts)
plt.show()
\end{minted}
\end{document}
